\section{Validasi Input di Server (PHP): Tipe Data, Format, dan Whitelist}

Validasi di server memastikan data yang masuk aman sebelum diproses. PHP menyediakan \code{filter\_var()} dengan filter bawaan: \code{FILTER\_VALIDATE\_EMAIL} (cek format email), \code{FILTER\_VALIDATE\_INT} (cek bilangan bulat), \code{FILTER\_VALIDATE\_URL} (cek format URL). Strategi \textbf{whitelist}: hanya terima nilai yang eksplisit diizinkan --- misalnya status hanya boleh ``aktif'' atau ``nonaktif''. Tolak input tidak valid dengan pesan error yang jelas tanpa mengungkap detail teknis \cite{phpRightWay}.

\begin{webcode}[language=PHP, caption={Validasi input lengkap di server}]
<?php
function sanitize(string $s): string {
  return htmlspecialchars(trim($s), ENT_QUOTES, 'UTF-8');
}

$errors = [];
$nama  = sanitize($_POST['nama'] ?? '');
$email = trim($_POST['email'] ?? '');
$umur  = $_POST['umur'] ?? '';
$role  = $_POST['role'] ?? '';

// Validasi panjang
if (strlen($nama) < 3 || strlen($nama) > 100) {
  $errors[] = "Nama harus 3-100 karakter.";
}
// Validasi format
if (!filter_var($email, FILTER_VALIDATE_EMAIL)) {
  $errors[] = "Format email tidak valid.";
}
// Validasi tipe dan rentang
$umur = filter_var($umur, FILTER_VALIDATE_INT,
  ['options' => ['min_range' => 17, 'max_range' => 100]]);
if ($umur === false) $errors[] = "Umur harus 17-100.";
// Whitelist
if (!in_array($role, ['mahasiswa', 'dosen', 'admin'])) {
  $errors[] = "Role tidak valid.";
}
?>
\end{webcode}

